\specsection{Measures on spaces of sequences}

One of the most highly developed branches of classical Probability Theory is that of limit theorems (law of large numbers, tendency towards the Gauss--Laplace law, \ldots); this concerns a deepening of the concept of statistical regularity manifested by phenomena that bring into play very large populations. The correct mathematical formulation of these problems requires the introduction of measures on sequence spaces; these spaces, which constitute the most obvious generalization of finite-dimensional spaces, were the subject of choice of investigations into `General Analysis' undertaken around 1920 by Fréchet, Lévy, Lusin, \ldots{} Nor is it fortuitous that Khintchine and Kolmogoroff, the creators of the new methods of Probability Theory, were both disciples of Lusin, and that Lévy very quickly oriented himself towards probabilistic problems: these constituted the touchstone of the new methods.

The first implicit intervention of a measure on a sequence space appeared in the work devoted by E. Borel in 1909 to countable probabilities (IV). A highly original idea of Borel consists in the application of the probabilistic results he had just obtained to the proof of properties possessed by the decimal expansion of almost every real number between 0 and 1. This application rests on the following fundamental remark: let us define every real number between 0 and 1 by the sequence of figures (or `digits') in its expansion in a given base \(q\) ( \(q \geqslant 2\) ); if one successively draws at random the various figures of a number \(x\) , independently of each other and with equal probability \(1 / q\) for \(0,1,\ldots ,q - 1\) , the probability that \(x\) lies in an interval of {[}0,1{[} is equal to the length of that interval.

In 1923, Steinhaus (V) established these results rigorously and described the precise mathematical model for the infinite sequence of random drawings considered by Borel: for simplicity let us take q=2 and denote by I the set with two elements \(\{0,1\}\) ; one equips I with the measure \(\mu\) defined by \(\mu(0)=\mu(1)=\frac{1}{2}\) ; the elements of the product space \(I^{N}\) are the sequences \(\varepsilon=(\varepsilon(n))_{n\in\mathbf{N}}\) of numbers equal to 0 or 1, and the mapping \(\varphi:\varepsilon\mapsto\sum_{n\geqslant0}\varepsilon(n)\cdot2^{-n-1}\) is, up to a countable set, a bijection of \(I^{N}\) onto the interval {[}0,1{]}; moreover, \(\varphi^{-1}\) transforms the Lebesgue measure on {[}0,1{]} into the measure P on \(I^{N}\) that is the product of the measures \(\mu\) on each of the factors. Actually, Steinhaus did not have at his disposal a construction for product measures; he used the existence of the quasi-bijection \(\varphi\) to construct the measure P on \(I^{N}\) starting from Lebesgue measure on {[}0,1{]}, and then gave an axiomatic characterization of P. The isomorphism so obtained made it possible to translate the language of probability into that of measure and to apply the known theorems on the Lebesgue integral.

In the same work, Steinhaus considered the random series \(\sum_{n\geqslant 0}\sigma_n\cdot a_n\) where the signs \(\sigma_{n} = \pm 1\) are chosen at random independently of each other and with equal probability \(\frac{1}{2}\) ; between 1928 and 1935, he studied numerous other random series. From their side, Paley, Wiener and Zygmund considered random Fourier series(1) of the form \(\sum_{n = -\infty}^{\infty}a_n\exp (2\pi i(nt + \Phi_n))\) ; the `amplitudes' \(a_{n}\) are fixed, and the `phases' \(\Phi_{n}\) are independent random variables uniformly distributed on {[}0,1{]}. While the analytic difficulties vary enormously from one of these problems to another, the translation in terms of measure theory is the same in all cases and represents an extension of the case treated by Borel and Steinhaus; it is a matter of constructing a measure on \(\mathbf{R}^{\mathbf{N}}\) that is the product of a family of measures all identical to a same positive measure \(\mu\) on \(\mathbf{R}\) of mass 1; for example, the preceding random Fourier series correspond to the case that \(\mu\) is the Lebesgue measure on {[}0,1{]}.

To construct such product measures, two methods can be used. The first is a direct method, precisely exposed for the first time by Daniell (VI, a)) in 1918; it was rediscovered in 1934 by Jessen (VII), who made a detailed study of the case that \(\mu\) is Lebesgue measure on {[}0, 1{]}. The second method is to seek artifices analogous to that of Steinhaus to reduce to Lebesgue measure on {[}0, 1{]}; this way of proceeding had the advantage of convenience so long as one did not have available a complete exposition of general measure theory, for it permitted using Lebesgue's theorems without the need for new proofs. \(^{(2)}\)

